\documentclass[../../uem_phy_std_q.tex]{subfiles}

\begin{document}
    \section{回路の諸々}
    \subsection{金属電子論}
        \subfile{01_metal-electron-theory/uem_phy_std_em_cx_01_q.tex}
        \vspace{20pt}
        \vfill
    \subsection{非オーム抵抗}
        \subfile{02_non-ohmic-resistor/uem_phy_std_em_cx_02_q.tex}
        \vspace{20pt}
        \vfill
    \subsection{合成抵抗・回路の対称性}
        \subfile{03_equivalent-resistance/uem_phy_std_em_cx_03_q.tex}
        \vspace{20pt}
        \vfill
    \subsection{コンデンサの無限回つなぎかえ①}
        \subfile{04_infinite-switching-1/uem_phy_std_em_cx_04_q.tex}
        \vspace{20pt}
        \vfill 
    \subsection{コンデンサの無限回つなぎかえ②}
        \subfile{05_infinite-switching-2/uem_phy_std_em_cx_05_q.tex}   
        \vspace{20pt}
        \vfill
    \subsection{誘電体}
        \subfile{06_dielectric/uem_phy_std_em_cx_06_q.tex}
        \vspace{20pt}
        \vfill
    \subsection{※導線中の自由電子の運動}
        \subfile{07_add-free-electrons/uem_phy_std_em_cx_07_q.tex}
        \vspace{20pt}
        \vfill
    \subsection{※電球を含む回路}
        \subfile{08_add-bulb/uem_phy_std_em_cx_08_q.tex}
        \vspace{20pt}
        \vfill
    \subsection{※合成抵抗}
        \subfile{09_add-equivalent-resistance/uem_phy_std_em_cx_09_q.tex}
        \vspace{20pt}
        \vfill
    \subsection{※誘電体を挿入したコンデンサ}
        \subfile{10_add-dielectric/uem_phy_std_em_cx_10_q.tex}
        \vspace{20pt}
        \vfill
\end{document}
